\section{4/1 - Compiler Optimization}
In the previous lecture, we tested that the compiler was valid. Optimizing the compiler is similar.
\\
A compiler has multiple phases:
\begin{enumerate}
  \item parse
  \item desugar
  \item typechecking/inference
  \item uniquify
  \item lowering
\end{enumerate}
But we could add in an optimization step that has a self loop
\begin{enumerate}
  \item parse
  \item desugar
  \item typechecking/inference
  \item uniquify
  \item \textbf{codegen}
  \item lowering
\end{enumerate}
OCaml uses Superoptimization, a pass to turn a program into the smallest possible representation.
\\
We plan to implement a compiler pass on the level of Calc in order to make our tree parsing more efficient. This code is built upon the languages from previous lectures.
\\
In our CalcLang we have expressions that look like this:
\begin{lstlisting}
let x = 1 + 2 in x
let ex0 = Let("x", App(Add, [Val (Nat 1); Val (Nat 2)]), Var "x")
\end{lstlisting}
One way we could optimize this is simply by performing the substitution and evaluation to get $3$. We can start by reducing our statically identifiable operations.
\begin{lstlisting}
let constfold_op (op, args) : CalcLang.expr =
  match op, args with
  | (Add, [Val (Nat n); Val (Nat m)]) -> Val (Nat (n + m))
  | (Add, [Val (Bool n); Val (Bool m)]) -> failwith "UNREACHABLE"
\end{lstlisting}
This works for this first simple case but it falls apart for the second one with a mistake. Though our typecheck should stop this but in the interest of simplicity we should learn how to deal with it seperately.
\newpage
Instead of operating over expressions, we might want to pass in a function.
\begin{lstlisting}
let constfold (e : expr) : expr =
  match e with 
  | Val v ->
  | Var x ->
  | Let (x, e1, e2) -> Let (x, constfold e1, constfold e2)
  | App (op, args) ->
    (match op, args with
    | (Add, [Val (Nat n); Val (Nat m)]) -> Val (Nat (n + m))
    | (Mul, [Val (Nat n); Val (Nat m)]) -> Val (Nat (n * m))
    | (Or, [Val (Bool n); Val (Bool m)]) -> Val (Nat (n || m))
    (* simply follow the eval op for the rest *)
    | _ -> App (op, List.map constfold args))
  | ITE (g, t, f) -> ITE(constfold g, constfold t, constfold f)
\end{lstlisting}
This code only works in the case where there is only one operation. With compound operations (such as 1 + (2 * 3)) it will only apply once.
\\
We want to continuously optimize our expression until we reach a fixed point.
\begin{lstlisting}
let rec optimize (opt : expr -> expr) (e : expr) : expr = 
  let e' = opt e in
  if e' = e
  then e
  else optimize opt e'
\end{lstlisting}
Once more we are faced with the problem of correctness. We could pass a trivial function (i.e. f(x)=0) but this doesn't preserve meaning.
\\
We could define a property to state our pass is valid $\forall e : Calc, e = constfold e$.
\\
We can also do some lazy boolean analysis:
\begin{lstlisting}
  | ITE (Val (Bool true), t, f) -> constfold t
  | ITE (Val (Bool false), t, f) -> constfold f
\end{lstlisting}
Now to optimize substitution. In a let binder we are looking for a case that looks like $e'[x \mapsto e]$.
\begin{lstlisting}
let rec substitute (x : string) (e : expr) (e' : expr) : expr =
  match e with
  | Val v -> Val v
  | Var x' -> if x' = x then e else Var x'
  | Let (x, e1, e2) -> 
    if x = x' then Let (x, e1, e2)
    else Let (x', substitute x e1, substitute x e2)
  | App (op, args) -> App(op, List.map (fun arg -> substitute x e arg))
  | ITE (g, t, f) -> ITE(substitute x e g, substitute x e t, substitute x e f)
\end{lstlisting}
We can now change our constfold function to use this in the Let case:
\begin{lstlisting}
  | Let (x, Val v, e2) -> substitute c (Val v) e2
  | Let (x, e1, e2) -> Let (x, constfold e1, constfold e2)
\end{lstlisting}
